\documentclass[10pt]{article}
\usepackage[margin=20mm]{geometry}
\usepackage{amsmath,amssymb,fontspec,unicode-math,xurl,hyperref}
\IfFontExistsTF{DejaVu Serif}{\setmainfont{DejaVu Serif}}{\setmainfont{Latin Modern Roman}}
\setmathfont{latinmodern-math.otf}
\setlength{\parindent}{0pt}\setlength{\parskip}{6pt}\sloppy
\begin{document}
Abstract: A treatment-dependent censoring example satisfies the stated identification assumptions but contradicts the explicitly defined weighting formula in Theorem 3.4. The discrepancy affects a core identification result, supporting Major impact within that stated domain. An arm-conditional correction is proved. Existing simulation and clinical estimates are not reproduced or refuted.\par
Author: AI-assisted technical review\par
Date:\par
Title: Technical Verification: Doubly Protected Estimation for Survival Outcomes Utilizing External Controls for Randomized Clinical Trials\par
\# Identity and scope\par
Doubly Protected Estimation for Survival Outcomes Utilizing External Controls for Randomized Clinical Trials. Chenyin Gao; Shu Yang; Mingyang Shan; Wenyu Ye; Ilya Lipkovich; Douglas Faries. PMLR 267 (2025), paper I25-43681. The reviewed source is the published PMLR PDF frozen for B004. Page numbers below are one-based PDF pages (also the article page labels), not volume pagination. This is an initial technical review; scientific adjudication remains pending.\par
The main paper, pp. 1–9, and mathematical Appendix A, pp. 12–25, were read, with a focused check of identification and its censoring assumptions. Appendix B and Figures 1–4 were inspected as reported simulations, not reproduced. Static inspection covered the released survival-estimation routine at the lines recorded in the coverage ledger. The study-specific clinical conclusions and original patient data were not independently validated.\par
\# Definitions and conditions\par
The target is the trial-population restricted mean survival-time (RMST) difference \(\theta_\tau=\int_0^\tau(S_1(t\mid R=1)-S_0(t\mid R=1))\,dt\). The observed data contain \(Y=\min(T,C)\) and \(\Delta=1(T<C)\); \(R=1\) denotes trial enrollment and \(A\) treatment. Assumption 3.1 requires internal treatment exchangeability and positive treatment/source probabilities. Assumption 3.2 requires \(T^{(a)}\perp C\mid X,A=a,R=r\). It does not require censoring to have the same distribution in the two randomized arms. Assumption 3.3 requires equality of control survival curves across sources, conditional on \(X\).\par
\# Substantive finding\par
Problem 1 (The stated identification formula omits treatment in the censoring probability).  Theorem 3.4(a), p. 3, defines \(\pi_1^C(t,X)=P(C\ge t\mid X,R=1)\) and uses this same denominator for both treatment arms. Under its stated Assumptions 3.1–3.3 this formula is false. This is a Major concern about a core identification result as stated; it does not establish failure in the narrower simulation designs.\par
\#\# Exact counterexample\par
Take \(X\) constant, \(P(R=1)=1/2\), and randomized \(P(A=1\mid R=1)=1/2\). In the external study \(A=0\). Let both potential event times have the exponential distribution with rate one, independently of assignment and source. Given \(R=1,A=0\), let \(C\sim\mathrm{Exp}(1)\); given \(R=1,A=1\), let \(C\sim\mathrm{Exp}(3)\). External censoring may be \(\mathrm{Exp}(1)\). In every arm draw censoring independently of the potential event times. Set, for example, \(\tau=1\) and check \(t=1/2\).\par
Consistency holds by construction. Treatment exchangeability holds because assignment is independent of potential outcomes. Both positivity probabilities equal \(1/2\). Conditional independent censoring holds in every observed arm. Control event-time distributions agree across sources for every \(t\). There are no ties, and every relevant censoring probability is positive at finite times. Theorem 3.4 does not invoke the additional asymptotic rate conditions of Assumption A.1.\par
Give the theorem the favorable interpretation that its treatment propensity is conditional on trial participation. Its marginal trial censoring survival is \[G(y)=\tfrac12e^{-y}+\tfrac12e^{-3y}.\] After cancellation of the source and treatment weights, the proposed formula for \(S_1(t\mid R=1)\) equals \[\begin{aligned}
I(t)&=\int_t^\infty\frac{e^{-y}e^{-3y}}{(e^{-y}+e^{-3y})/2}\,dy\\
&=2\int_0^{e^{-t}}\frac{u^2}{1+u^2}\,du
=2\{e^{-t}-\arctan(e^{-t})\}.\end{aligned}\] The correct survival is \(e^{-t}\). At \(t=1/2\), the claimed expression is \(0.1226460717641\) and the true value is \(0.6065306597126\). Their difference is a population-level discrepancy, with no estimation or finite-sample nuisance error. The independent script checks the integral and closed form with SciPy quadrature; both agree to the reported digits. Even at \(t=0\) the proposed expression is \(2-\pi/2\), not one.\par
\#\# Counterevidence and narrower interpretation\par
Appendix A.1, p. 12, ends with a denominator \(S^C(Y\mid X,A=0,R=1)\) for the control arm. This is the appropriately arm-conditional quantity, and is material evidence that an omitted conditioning variable may have been intended in the main theorem. It is not the explicitly defined marginal quantity on p. 3. The appendix also conditions \(T^{(0)}>t\) on \(C>T\) in an intermediate equality; independent event and censoring times do not justify conditioning on the event \(C>T\), which selects longer or shorter event times unequally. The corrected weighting identity below avoids that step.\par
The simulation censoring hazard in Section 4 depends on \(X\) and the study but not treatment. Those experiments lie in a special case where the present distinction vanishes. Static code inspection found that `srEC\_Surv.R`, lines 27–34 and 62–71, defaults `X.pred.c` to the three \(X\) variables and uses that list in the censoring model, while `X.pred.f` includes \(A\) and is used in the event model. The argument documentation names those two lists in the opposite order. Thus defaults do not silently impose an arm-specific censoring fit, though a caller could explicitly add \(A\) to the list actually used for censoring. No claim is made about the exact call used for historical results.\par
\#\# Verified repair of the identification identity\par
Define \(G_{a,1}(y,X)=P(C\ge y\mid X,A=a,R=1)\) and use the trial treatment propensity \(p_a(X)=P(A=a\mid X,R=1)\). Then \[S_a(t\mid R=1)=\frac1{P(R=1)}\mathbb E\!\left[\frac{R1(A=a)\Delta1(Y>t)}{p_a(X)G_{a,1}(Y,X)}\right]\] under the stated conditional independence, consistency, treatment exchangeability and positive denominators. Conditional on \(T,X,A=a,R=1\), the expectation of \(1(C>T)/G_{a,1}(T,X)\) is one (for continuous times). Iterated expectation then gives the result. In the example this replaces \(G(y)\) by \(e^{-3y}\) and yields \(\int_t^\infty e^{-y}\,dy=e^{-t}\), independently verified numerically. Alternatively, adding \(C\perp A\mid X,R=1\) restores the marginal denominator. These are different repairs: the former changes the formula, the latter narrows the assumptions. This report does not certify all efficiency projections or asymptotic claims after either change.\par
\# Other checked reasoning and limitations\par
The source/source-odds identity in Theorem 3.4(b) follows by iterated expectation for the external controls on common support once appropriate censoring weights are used. Theorem 3.5 then uses survival-curve identification as the foundation for the integrative influence functions, establishing the core dependency of Problem 1. The nuisance-product expression in Theorem 3.6 must be interpreted as a remainder bound rather than a bound on the entire random estimation error: even known nuisances leave the empirical influence-function term. Appendix A.3 mixes these objects, so this review does not certify its exact rate statement.\par
The selective procedure on pp. 5–6 detects zero integrated RMST drift whereas Appendix A.5 invokes equality of the survival curves for all times in the selected set. These are distinct restrictions. As a bounded diagnostic, \(S_T(t)=e^{-t}\) and \(S_E(t)=a e^{-t/2}+(1-a)e^{-2t}\) have equal RMST at \(\tau=1\) for \[a=\frac{(1-e^{-1})-(1-e^{-2})/2}{2(1-e^{-1/2})-(1-e^{-2})/2}=0.5634084559\ldots,\] while the curves differ. This does not satisfy Assumption 3.3 and is not offered as a counterexample to a theorem requiring that assumption. It identifies remaining work in extending the adaptive claim beyond homogeneous data: zero RMST drift alone cannot establish the curve-level tangent-space restriction. The optional oracle-weighted-curve calculation in the evidence is only a finite diagnostic of this distinction, not an empirical reproduction or a second adjudicated theorem counterexample.\par
Figures 1 and 3 show simulation trends for the specified settings, with some finite-sample type-I-error and efficiency limitations acknowledged in the Discussion. Figure 2 is an application analysis rather than independent evidence of transportability assumptions. None of these plots tests the arm-dependent-censoring counterexample. The original clinical estimates can survive this concern; determining their sensitivity requires actual censoring models, call configurations and patient-level data. The highest supported assessment concerns the false identification formula under its stated domain. No later erratum was established from the frozen packet.\par
\# Recorded state\par
Dimension | State\par
Inspection assessment | Major\par
Scientific adjudication | Pending\par
Repair verification | Limited to statements explicitly proved here\par
Author adoption | Not established\par
Human reviews | None supplied\par
9 Chenyin Gao; Shu Yang; Mingyang Shan; Wenyu Ye; Ilya Lipkovich; Douglas Faries. Doubly Protected Estimation for Survival Outcomes Utilizing External Controls for Randomized Clinical Trials. PMLR 267, 2025. Frozen published version. https://proceedings.mlr.press/v267/gao25s.html. Released source: https://github.com/Gaochenyin/SelectiveIntegrative/tree/c126485dda959982198b3deb93487e1a0a24350f. Snapshot inspected only in the files and lines listed in coverage.json; historical experimental provenance is unverified.\par
\# Evidence and reproducibility\par
The source SHA-256 is `7e0c643099b045f65c49b8b32049b93beaf9b9a37a9bc600a341f50ee5b12d1c`. Independent calculations, their exact outputs, input hashes, static code anchors, coverage and build records are stored under `verification/0312/`. The script `independent\_check.py` is reviewer-written and imports no author program. Numerical checks support the explicit bounded arguments above and are not benchmark reproductions. Source illustrations inspected as images are identified in the coverage ledger. All report pages are subject to a separate hash-bound visual record after compilation. AI performed this initial review; no human review or independent acceptance is claimed. Separate OpenReview attachments inaccessible during source preparation remain unavailable. Later revisions and errata beyond the frozen source packet were not independently searched; no novelty or priority of discovery is asserted.\par
\end{document}
